 % ------------ ELM-proceedings-template.tex  -- ELM proceedings template ---------------------------------------------------------------------------
% Adapted by Alex Kalomoiros and Florian Schwarz from the  LSA Proceedings LaTeX template created by Sarah E. Murray.
% http://journals.linguisticsociety.org/proceedings/index.php/PLSA/pages/view/instructions
%
%
% ------------ begin preamble -----------------------------------------------------------------------------------------
 \documentclass[12pt,letterpaper]{article}		

% ------------ include additional packages for your own use here. Check Preamble.tex for what packages are already included, and to resolve potential conflicts between packages.

%----------------------------------------------------------------------------

\include{Preamble}

% -------------- Place personal definitions of new commands etc. here ---------------------------------------------------------------------------------------------------------   


%-------------------------------------------------------------- title information --------------------------------------%
\title{Style sheet and \LaTeX template for the \textit{Proceedings of ELM}}
%
\author{First Last, First Last, \& First Last\footnote{Both acknowledgments and author affiliation information go in an initial footnote like this, referenced by an asterisk on the end of the last author's name on the author-name line at the top. You can say who you would like to thank here and then end the footnote as follows. Authors: Author One, University of the Atlantic (\href{mailto:author1@atlantic.edu}{author1@atlantic.edu}) \& Other Author, Pacific University (\href{mailto:author2@gmail.com}{author2@gmail.com}).}}


% ------------ end preamble ------------------------------------------------------------------------------------------------

% ------------ begin main document --------------------------------------------------------------------------------------
\begin{document} 

%%If using linguex, need the following commands to get correct LSA style spacing
%% these have to be after  \begin{document}
\setlength{\Extopsep}{6pt}
\setlength{\Exlabelsep}{6.75pt}		%effect of 0.3in indent from left text edge
%%
 
\maketitle
\thispagestyle{empty}
\begin{abstract}
This is the \LaTeX template (and/or the pdf generated from it) for the Proceedings of ELM, published by the LSA. (A Word version with matching formatting setup and instructions is available here: \url{https://elm-conference.net/wp-content/proceedings/ELM_Word_Template.docx}.) The title and author name(s) paragraphs, centered at the top of the first page, and the abstract and keyword paragraphs should match what is entered for these fields in the submission procedure. These will all be included on the landing page for the article and in the metadata that accompanies it for bibliographical and search-engine indexing purposes. The editorial team will add citation information with DOI in the header area of the first page and a copyright notice in the footer showing the year of publication and that the author or authors own the copyright and permit publication by the Linguistic Society of America under a CC BY license. The editorial team will also handle the page numbering details at the final layout stage.
\end{abstract}

\begin{keywords} %separated by semicolons
ELM; semantics; pragmatics; psycholinguistics; proceedings; templates; style sheets; formatting; LSA uniform style
\end{keywords}

\vspace{6pt} % Needed for proper spacing betwen Keywords and first paragraph.


\section{Introduction} %Section titles will be displayed inline
Your article needs to have the formatting, layout, and style features displayed in this document. This is the \LaTeX template adapted for ELM from the LSA Proceedings stylesheet.\footnote{\url{https://drive.google.com/file/d/1wOQJkLJHOI2WTWYjPmb23hZIoit0IWOz/view}} (A Word version with matching formatting setup and instructions is available here: \url{https://www.elm-conference.net/wp-content/uploads/2024/07/ELM_Word_Template.docx}.) The required features, already fully in place in this document, include the following:

\begin{itemize}
\item	paper: 8.5 x 11 inches with margins of 1 inch at the top, bottom, left, and right;
\item	font family: Times New Roman or Times;
\item	font size: 12 points by default; 10 points for footnotes; font property = small caps for subsection titles and interlinear glosses;
\item	line height: single-spacing;
\item	text alignment: justified (hyphenation allowed);
\item	spacing between sections: 6 points by default; 12 points between the abstract and the author name line, between the first paragraph of the body and the keywords, and between the last paragraph and the list of references.
\end{itemize}

Note that section titles for content in the body (i.e., after the keywords paragraph) are in bold, are numbered, and are wholly contained within the first paragraph of the section that they name. These and all titles use sentence case, which is to say that only the initial ordinary word or any ordinary word after a colon has the first letter in uppercase. Naturally, words that ordinarily have the first letter in uppercase, such as proper names, should have them that way in sentence case as well.

The non-initial paragraphs of body text within each section should be indented by 0.3 inches, like this one, and have no extra space above them, like this one.
 
 
 
 \section{Further Details} 
This section, which like all sections but the first, has a space of 6 points between it and the preceding paragraph, illustrates how to include subsections, if you choose to use them, and how to include language examples and illustrations, tables, figures, and the like. If you choose to use subsections, you should probably not go deeper than one to two levels of subsections. Subsections are numbered (in regular font) like sections, and differ only in their font, which uses small caps.

\subsection{How to include language examples}

 For interlinear glosses use the abbreviation system and formatting conventions of Leipzig Glossing Rules,\footnote{\url{http://www.eva.mpg.de/lingua/resources/glossing-rules.php}} but do not italicize any of the lines or indent the number labels. The examples in this document use the linguex package and have been adjusted to the LSA format requirements (see commands right after \verb=\begin{document}= in the \LaTeX code). For glossed examples, append `g' to the linguex example commands (e..g, \verb=\exg.= or \verb=\ag.=; see linguex documentation for details.\footnote{\url{http://texdoc.net/texmf-dist/doc/latex/linguex/linguex-doc.pdf}}

Example \ref{label} below shows how to include a single numbered language example with interlinear glossing together with a line identifying the language and its source.

\exg. Gila	abur-u-n	ferma	hamišaluǧ g\"uǧ\"una	amuq’-da-č.\\
Gila	they-\Obl{}-\Gen{}	farm	forever	behind	stay-\Fut{}-\Neg{}\\
`Now their farm will not stay behind forever.' \hfill (Lezgian, \citet{haspelmath93}:207)\footnote{The Lezgian alphabet is used here, as in the source document. IPA symbols need not be used when presenting examples in the manner used in a source document or from languages with their own alphabet. But, the characters should be explained with reference to IPA characters, in a footnote, for example (Lezgian \v{s} =  /\textipa{S}/, v{g} = /g$^w$/, etc.).}\label{label}


Example  \ref{label2} shows how to include multiple letter-labeled examples under a single number label.

\ex. \label{label2}
\ag. \~{n}uka-ka	Maria papa-ta	yanu-chi-shka	ka-rka-ni.\\
\Fsg{}-\Top{}	Maria	potato-\Acc{}	cook-\Caus{}-\Pass{}	be-\Pst{}-\Fsg{}.\leipzigfont{sbg} \\
`I was made to cook potatoes by Maria.'
\bg. *ñuka-ka	mishqui-ka	miku-naya-shka	ka-rka-nil.\\
\Fsg{}-\Top{}	candy-\Top{}	eat-\leipzigfont{desid}-\Pass{}	be-\Pst{}-\Fsg{}.\leipzigfont{sbg}\\
`I wanted to eat candy.' 

There should be no extra space between any contiguous examples. But each whole set of such should have a 6-point space above and below. The first line may be set in italics, if you prefer.

\subsection{Figures, drawings, and tables}

Center tables, figures, and other illustrative items and include a centered caption line beneath them. Provide 6 points of vertical space above and below the entire array and between the caption and the item. Figures and tables must be referenced initially in the body text at some point before they appear. And, they should appear as close to their initial reference as is feasible. An illustration is provided in Figure 1, using the graphicx package. All included figures should follow the format below.\footnote{This example picture is taken from \citet{beltrschw}.}

\begin{figure}[ht]
\centering
    \includegraphics[scale = 0.4]{Example}\\
    \caption{Choices by participants' identity: Prolific \label{f:figure}}
\end{figure}

An example of table formatting is illustrated by Table \ref{t:table}. It uses the same formatting as Figure \ref{f:figure}, but with Table \textit{N} as the caption label. Use only simple horizontal lines as separators, if desired. The spacing before and after this table is controlled by the \verb=\setlength\intextsep{6pt}= command in the preamble. The space between the table and the caption is also set in the preamble, by \verb=\setlength\abovecaptionskip{6pt}=.
 

\begin{table}[ht]
\centering
	\begin{tabular}{l c c}
	& Heading 1 & Heading 2\\
	\hline
	X &  Some data & Some data\\
	Y & Some data & Some data \\
	\hline
	\end{tabular}
\caption{Table showing some data about something \label{t:table}}
\end{table}



Drawings of syntactic trees should be added as numbered examples, parallel to other linguistic examples. Authors can choose on their own which tree generating packages to use in \LaTeX. 

 \section{Supplementary materials} If you have supplementary materials that you want to make avail- able to readers, such as video clips, data files, or large images, and they cannot be accommodated in a PDF article of this kind, you can ask the editors to upload these to a dedicated proceedings resources space on the ELM website, or host them on a personal or institutional external website or on such media-sharing sites as YouTube, Dropbox, or Google Drive, and provide active hyperlinks to them in the body of your article, using the \verb=\url{URL}= command. (Note that links introduced with \verb=\href= will not be preserved in the final 

\section{Citations and references} \label{s:cite} Use author-date notation for in-text citations, for example: \ldots as noted recently by \citet{Silk2012} and \citet{AnderBois2014} , drawing on insights from various other researchers \cite[e.g.][]{Chomsky1986,Yu2003,Christiansenetal2016}, there have been numerous technological advances in the procedures used to publish research articles online. One of these is the inclusion of URLs or DOIs with active hyperlinks in lists of references, as illustrated below. Include DOIs or URLs, when available. For all other matters, follow the formatting conventions of the LSA uniform style sheet\footnote{\url{http://www.linguisticsociety.org/sites/default/files/USS-NoComments.pdf}} for the list of references. 

The Semantics \& Pragmatics stylesheet, sp.bst\footnote{ \url{https://raw.githubusercontent.com/semprag/tex/master/sp.bst}} meets the LSA formatting requirements, and so is used below. See \url{http://info.semprag.org} for further details. 

A note on citing conference proceedings, from the Unified style sheet for linguistics (\href{https://www.linguisticsociety.org/sites/default/files/style-sheet_0.pdf}):

\begin{quoting}
    ``For conference proceedings
published with an ISSN, treat the proceedings as a journal: Include both the full
conference name and any commonly used acronym for the conference (BLS, WCCFL,
etc.) in the journal title position. For proceedings not published with an ISSN, treat the
proceedings as any other book, using the full title as listed on the front cover or title
page. If the title (and subtitle if there is one) only includes an acronym for the
conference name, expand the acronym in square brackets or parentheses following the
acronym. If the title does not include an acronym which is commonly used for the
conference name, include the acronym in square brackets or parentheses following the
conference name. The advantage of including the acronym after the society title is that it
makes the entry much more identifiable in a list of references.''

\end{quoting}

\noindent To accomplish the correct formatting using BibTeX, simply code the bib-entries as \verb=@article= or \verb=@incollection= accordingly.
 
% ------------ references --------------------------------------------------------------------------------------
\setlength{\bibhang}{0.3in}			% hanging indent for references must be 0.3in
\titleformat{\section}{\normalfont\bfseries}{\thesection}{.5em}{}		% references section not supposed to be followed by a period


\bibliographystyle{sp}		% S\&P bibliography style
%S\&P, an LSA publication, meets the LSA guidelines. 
%See: http://info.semprag.org 
%and get the bst at: https://raw.githubusercontent.com/semprag/tex/master/sp.bst 
%If using the sp.bst, the following command turns your DOIs into links
\newcommand{\doi}[1]{\href{#1}{#1}}	%modified from sp.cls
\vspace{6pt}

\bibliography{sample}			% CHANGE 'SAMPLE' TO YOUR BIB FILE AS NECESSARY


\end{document}
 
 
 
 